%% ============================================================
%%  SECCIÓN IX — Tratamiento de Datos y Resultados
%% ============================================================

\section{Tratamiento de Datos y Resultados}
\label{sec:tratamiento}

En esta sección se obtienen, a partir de los datos experimentales de la
Sección~\ref{sec:resultados}, las magnitudes derivadas necesarias para
caracterizar el comportamiento elástico de cada material. Se adopta
$g = \qty{9{,}8}{m/s^2}$ y se emplean exclusivamente las masas reales
registradas con la balanza.

%% ------------------------------------------------------------------
\subsection{Área de la sección transversal inicial}
%% ------------------------------------------------------------------

\subsubsection{Resorte metálico — sección circular}

\begin{equation}
    S_{0,r} = \frac{\pi\,D_0^2}{4}
    \label{eq:area_resorte}
\end{equation}

Sustituyendo $D_0 = \qty{14{,}00}{mm} = \qty{1{,}400e-2}{m}$:

\begin{equation*}
    S_{0,r} = \frac{\pi\,(1{,}400\times10^{-2})^2}{4}
             = \qty{1{,}539e-4}{m^2}
\end{equation*}

\subsubsection{Liga elástica — sección rectangular}

\begin{equation}
    S_{0,l} = a_0\,b_0
    \label{eq:area_liga}
\end{equation}

Sustituyendo $a_0 = \qty{14{,}80}{mm}$ y $b_0 = \qty{5{,}50}{mm}$:

\begin{equation*}
    S_{0,l} = (14{,}80\times10^{-3})(5{,}50\times10^{-3})
             = \qty{8{,}14e-5}{m^2}
\end{equation*}

%% ------------------------------------------------------------------
\subsection{Magnitudes derivadas — Resorte metálico}
%% -----------------------------------------------------------------

Para cada configuración de carga $i$ se calculan el área deformada $S_i$ y el esfuerzo real $\sigma_i$:

\begin{equation}
    F_i = m_i\,g, \qquad
    \Delta\ell_i = L_i - L_0, \qquad
    \varepsilon_i = \frac{\Delta\ell_i}{L_0}, \qquad
    S_i = \frac{\pi\,D_i^2}{4}, \qquad
    \sigma_i = \frac{F_i}{S_i}
    \label{eq:mag_deriv}
\end{equation}

con $L_0 = \qty{16{,}60}{cm} = \qty{0{,}1660}{m}$.

\paragraph{Cálculo representativo para $i = 1$:}

\begin{align*}
    F_1        &= 0{,}4998 \times 9{,}8 = \qty{4{,}898}{N}\\
    \Delta\ell_1 &= 0{,}298 - 0{,}1660  = \qty{0{,}1320}{m}\\
    \varepsilon_1 &= \frac{0{,}1320}{0{,}1660} = 0{,}7952\\
    S_1 &= \frac{\pi\,(14{,}22\times10^{-3})^2}{4} = \qty{1{,}588e-4}{m^2}\\
    \sigma_1  &= \frac{4{,}898}{1{,}588\times10^{-4}} = \qty{30{,}84}{kPa}
\end{align*}

\begin{table}[H]
    \centering
    \caption{Magnitudes derivadas del resorte metálico empleando esfuerzo real.}
    \label{tab:deriv_resorte}
    \begin{tabular}{cccccc}
        \toprule
        \textbf{N$^{\circ}$}
            & \textbf{$F$ (N)}
            & \textbf{$\Delta\ell$ (m)}
            & \textbf{$\varepsilon$}
            & \textbf{$S_i$ ($10^{-4}$ m$^2$)}
            & \textbf{$\sigma$ (kPa)} \\
        \midrule
        1 &  4{,}898 & 0{,}1320 & 0{,}7952 & 1{,}588 &  30{,}84 \\
        2 &  7{,}392 & 0{,}1550 & 0{,}9337 & 1{,}570 &  47{,}08 \\
        3 &  9{,}849 & 0{,}1840 & 1{,}1084 & 1{,}548 &  63{,}62 \\
        4 & 12{,}343 & 0{,}2190 & 1{,}3193 & 1{,}535 &  80{,}41 \\
        5 & 14{,}747 & 0{,}2510 & 1{,}5120 & 1{,}526 &  96{,}64 \\
        6 & 17{,}241 & 0{,}2800 & 1{,}6867 & 1{,}517 & 113{,}65 \\
        \bottomrule
    \end{tabular}
\end{table}

%% ------------------------------------------------------------------
\subsection{Magnitudes derivadas — Liga elástica}
%% ------------------------------------------------------------------

Se aplican relaciones análogas empleando $S_i = a_i \times b_i$ 
y $L_0 = \qty{34{,}20}{cm} = \qty{0{,}3420}{m}$.

\paragraph{Cálculo representativo para $i = 1$ (carga):}

\begin{align*}
    F_1          &= \qty{4{,}898}{N}\\
    \Delta\ell_1 &= 0{,}367 - 0{,}3420 = \qty{0{,}0250}{m}\\
    \varepsilon_1 &= \frac{0{,}0250}{0{,}3420} = 0{,}07310\\
    S_1 &= (14{,}24\times10^{-3})(5{,}60\times10^{-3}) = \qty{79{,}74e-6}{m^2}\\
    \sigma_1  &= \frac{4{,}898}{79{,}74\times10^{-6}} = \qty{61{,}42}{kPa}
\end{align*}

\begin{table}[H]
    \centering
    \caption{Magnitudes derivadas de la liga elástica — fase de carga.}
    \label{tab:deriv_liga_carga}
    \begin{tabular}{cccccc}
        \toprule
        \textbf{N$^{\circ}$}
            & \textbf{$F$ (N)}
            & \textbf{$\Delta\ell$ (m)}
            & \textbf{$\varepsilon$}
            & \textbf{$S_i$ ($10^{-6}$ m$^2$)}
            & \textbf{$\sigma$ (kPa)} \\
        \midrule
        1 &  4{,}898 & 0{,}0250 & 0{,}07310 & 79{,}74 &  61{,}42 \\
        2 &  7{,}392 & 0{,}0440 & 0{,}12865 & 75{,}13 &  98{,}39 \\
        3 &  9{,}849 & 0{,}0730 & 0{,}21345 & 68{,}85 & 143{,}05 \\
        4 & 12{,}343 & 0{,}1080 & 0{,}31579 & 68{,}85 & 179{,}27 \\
        5 & 14{,}747 & 0{,}1290 & 0{,}37719 & 67{,}44 & 218{,}67 \\
        6 & 17{,}241 & 0{,}1410 & 0{,}41228 & 66{,}00 & 261{,}23 \\
        \bottomrule
    \end{tabular}
\end{table}

\begin{table}[H]
    \centering
    \caption{Magnitudes derivadas de la liga elástica — fase de descarga.}
    \label{tab:deriv_liga_descarga}
    \begin{tabular}{cccccc}
        \toprule
        \textbf{N$^{\circ}$}
            & \textbf{$F$ (N)}
            & \textbf{$\Delta\ell$ (m)}
            & \textbf{$\varepsilon$}
            & \textbf{$S_i$ ($10^{-6}$ m$^2$)}
            & \textbf{$\sigma$ (kPa)} \\
        \midrule
        5 & 14{,}747 & 0{,}1320 & 0{,}38596 & 63{,}08 & 233{,}78 \\
        4 & 12{,}343 & 0{,}1180 & 0{,}34503 & 71{,}82 & 171{,}86 \\
        3 &  9{,}849 & 0{,}1130 & 0{,}33041 & 72{,}74 & 135{,}40 \\
        2 &  7{,}392 & 0{,}0800 & 0{,}23392 & 72{,}47 & 101{,}95 \\
        1 &  4{,}898 & 0{,}0520 & 0{,}15205 & 75{,}89 &  64{,}54 \\
        \bottomrule
    \end{tabular}
\end{table}

%% ------------------------------------------------------------------
\subsection{Constante recuperadora del resorte}
%% ------------------------------------------------------------------

Por la Ley de Hooke $F = k\,\Delta\ell$, la constante $k$ se obtiene
como la pendiente del ajuste lineal mínimo-cuadrático a los datos
$(\Delta\ell_i,\,F_i)$:

\begin{equation}
    k = \frac{n\displaystyle\sum_{i}\Delta\ell_i F_i
              - \displaystyle\sum_{i}\Delta\ell_i\displaystyle\sum_{i}F_i}
             {n\displaystyle\sum_{i}(\Delta\ell_i)^2
              - \!\left(\displaystyle\sum_{i}\Delta\ell_i\right)^{\!2}}
    \label{eq:k_lsq}
\end{equation}

\begin{table}[H]
    \centering
    \caption{Productos auxiliares para el ajuste $F$ vs.\ $\Delta\ell$ del resorte.}
    \label{tab:sumas_resorte}
    \begin{tabular}{crrr}
        \toprule
        \textbf{N$^{\circ}$}
            & \textbf{$\Delta\ell_i$ (m)}
            & \textbf{$(\Delta\ell_i)^2$ (m$^2$)}
            & \textbf{$\Delta\ell_i\,F_i$ (N$\cdot$m)} \\
        \midrule
        1 & 0{,}1320 & 0{,}017424 & 0{,}64654 \\
        2 & 0{,}1550 & 0{,}024025 & 1{,}14576 \\
        3 & 0{,}1840 & 0{,}033856 & 1{,}81222 \\
        4 & 0{,}2190 & 0{,}047961 & 2{,}70312 \\
        5 & 0{,}2510 & 0{,}063001 & 3{,}70150 \\
        6 & 0{,}2800 & 0{,}078400 & 4{,}82748 \\
        \midrule
        $\Sigma$ & 1{,}2210 & 0{,}264667 & 14{,}83662 \\
        \bottomrule
    \end{tabular}
\end{table}

Sustituyendo en~\eqref{eq:k_lsq} con $n = 6$ y
$\sum F_i = \qty{66{,}470}{N}$:

\begin{equation*}
    k = \frac{6(14{,}83662) - (1{,}2210)(66{,}470)}
             {6(0{,}264667) - (1{,}2210)^2}
      = \frac{89{,}020 - 81{,}160}
             {1{,}5880 - 1{,}4908}
      = \frac{7{,}860}{0{,}09718}
\end{equation*}

La incertidumbre de $k$ se estima mediante el error estándar de la
pendiente del ajuste por mínimos cuadrados, evaluando numéricamente
$s_k = \qty{2{,}38}{N/m}$, por lo que:

\begin{equation}
    \boxed{k = \qty{80{,}89 \pm 2{,}38}{N/m}}
    \label{eq:resultado_k}
\end{equation}

%% ------------------------------------------------------------------
\subsection{Módulo de Young}
%% ------------------------------------------------------------------

El módulo de Young $Y$ corresponde a la pendiente de la recta de ajuste
en el diagrama $\sigma$ vs.\ $\varepsilon$:

\begin{equation}
    Y = \frac{n\displaystyle\sum_{i}\varepsilon_i\,\sigma_{i}
              - \displaystyle\sum_{i}\varepsilon_i\displaystyle\sum_{i}\sigma_{i}}
             {n\displaystyle\sum_{i}\varepsilon_i^2
              - \!\left(\displaystyle\sum_{i}\varepsilon_i\right)^{\!2}}
    \label{eq:Young_lsq}
\end{equation}

\subsubsection{Resorte metálico}

\begin{table}[H]
    \centering
    \caption{Productos auxiliares para el ajuste $\sigma$ vs.\ $\varepsilon$
             del resorte.}
    \label{tab:sumas_young_resorte}
    \begin{tabular}{crrr}
        \toprule
        \textbf{N$^{\circ}$}
            & \textbf{$\varepsilon_i$}
            & \textbf{$\varepsilon_i^2$}
            & \textbf{$\varepsilon_i\,\sigma_i$ (kPa)} \\
        \midrule
        1 & 0{,}7952 & 0{,}63235 &  24{,}524 \\
        2 & 0{,}9337 & 0{,}87179 &  43{,}959 \\
        3 & 1{,}1084 & 1{,}22855 &  70{,}516 \\
        4 & 1{,}3193 & 1{,}74055 & 106{,}085 \\
        5 & 1{,}5120 & 2{,}28614 & 146{,}120 \\
        6 & 1{,}6867 & 2{,}84497 & 191{,}695 \\
        \midrule
        $\Sigma$ & 7{,}3553 & 9{,}60435 & 582{,}899 \\
        \bottomrule
    \end{tabular}
\end{table}

Con $\sum\sigma_{i} = \qty{432{,}24}{kPa}$ y $n = 6$:

\begin{equation*}
    Y_r = \frac{6(582{,}899) - (7{,}3553)(432{,}24)}
               {6(9{,}60435) - (7{,}3553)^2}
        = \frac{3497{,}394 - 3179{,}256}
               {57{,}6261 - 54{,}1004}
        = \frac{318{,}138}{3{,}5257}
\end{equation*}

Evaluando numéricamente $s_Y = \qty{2{,}65}{kPa}$,
por lo que:

\begin{equation}
    \boxed{Y_r = \qty{90{,}23 \pm 2{,}65}{kPa}}
    \label{eq:Young_resorte}
\end{equation}

\subsubsection{Liga elástica — fase de carga}

Se emplea únicamente la fase de carga para el ajuste lineal inicial.

\begin{table}[H]
    \centering
    \caption{Productos auxiliares para el ajuste $\sigma$ vs.\ $\varepsilon$
             de la liga (carga).}
    \label{tab:sumas_young_liga}
    \begin{tabular}{crrr}
        \toprule
        \textbf{N$^{\circ}$}
            & \textbf{$\varepsilon_i$}
            & \textbf{$\varepsilon_i^2$}
            & \textbf{$\varepsilon_i\,\sigma_i$ (kPa)} \\
        \midrule
        1 & 0{,}07310 & 0{,}005344 &   4{,}490 \\
        2 & 0{,}12865 & 0{,}016551 &  12{,}658 \\
        3 & 0{,}21345 & 0{,}045561 &  30{,}534 \\
        4 & 0{,}31579 & 0{,}099723 &  56{,}612 \\
        5 & 0{,}37719 & 0{,}142272 &  82{,}480 \\
        6 & 0{,}41228 & 0{,}169975 & 107{,}700 \\
        \midrule
        $\Sigma$ & 1{,}52046 & 0{,}479426 & 294{,}474 \\
        \bottomrule
    \end{tabular}
\end{table}

Con $\sum\sigma_i = \qty{962{,}03}{kPa}$ y $n = 6$:

\begin{equation*}
    Y_l = \frac{6(294{,}474) - (1{,}52046)(962{,}03)}
               {6(0{,}479426) - (1{,}52046)^2}
        = \frac{1766{,}844 - 1462{,}728}
               {2{,}87656 - 2{,}31180}
        = \frac{304{,}116}{0{,}56476}
\end{equation*}

\begin{equation}
    \boxed{Y_l = \qty{538 \pm 25}{kPa} \approx \qty{0{,}538}{MPa}}
    \label{eq:Young_liga}
\end{equation}

Dado el carácter no lineal de la liga, $Y_l$ debe interpretarse como un
módulo de Young efectivo o \textit{módulo secante} promediado sobre la
región de carga ensayada, no como una propiedad intrínseca del material.
El análisis riguroso requeriría el módulo tangente
$E_t = \mathrm{d}\sigma/\mathrm{d}\varepsilon$, dependiente de la
deformación.

%% ------------------------------------------------------------------
\subsection{Trabajo de deformación del resorte}
%% ------------------------------------------------------------------

El trabajo realizado se estima mediante la regla trapezoidal aplicada 
a la curva $F$ vs.\ $\Delta\ell$:

\begin{equation}
    W \approx \sum_{i=0}^{2}
      \frac{F_i + F_{i+1}}{2}\,(\Delta\ell_{i+1} - \Delta\ell_i)
    \label{eq:trabajo_trap}
\end{equation}

donde $(\Delta\ell_0,\,F_0) = (0,\,0)$ es el estado inicial sin carga.
Sustituyendo:

\begin{align*}
    W &\approx \frac{0 + 4{,}898}{2}(0{,}1320)
             + \frac{4{,}898 + 7{,}392}{2}(0{,}0230)
             + \frac{7{,}392 + 9{,}849}{2}(0{,}0290)\\[6pt]
      &= 0{,}3233 + 0{,}1413 + 0{,}2500
\end{align*}

\begin{equation}
    \boxed{W = \qty{0{,}715}{J}}
    \label{eq:trabajo}
\end{equation}

%% ------------------------------------------------------------------
\subsection{Área del lazo de histéresis}
%% ------------------------------------------------------------------

El área encerrada por el lazo de histéresis en el diagrama 
$\sigma$--$\varepsilon$ representa la energía disipada por unidad de 
volumen. Se calcula mediante la integral cerrada:

\begin{equation}
    A = \oint \sigma \, d\varepsilon 
      = \int_{\text{carga}} \sigma \, d\varepsilon 
      - \int_{\text{descarga}} \sigma \, d\varepsilon
    \label{eq:area_histeresis}
\end{equation}

Evaluando ambas integrales numéricamente mediante la regla trapezoidal
($\sum \frac{\sigma_i + \sigma_{i+1}}{2} \Delta\varepsilon$):

\begin{table}[H]
    \centering
    \caption{Integración numérica del área bajo las curvas de carga y descarga.}
    \label{tab:area_histeresis}
    \begin{tabular}{cc|cc}
        \toprule
        \multicolumn{2}{c|}{\textbf{Fase de Carga}} & \multicolumn{2}{c}{\textbf{Fase de Descarga}} \\
        \textbf{Intervalo} & \textbf{Área ($kJ/m^3$)} & \textbf{Intervalo} & \textbf{Área ($kJ/m^3$)} \\
        \midrule
        0 $\to$ 1 &  2{,}245 & 6 $\to$ 5d &  6{,}514 \\
        1 $\to$ 2 &  4{,}439 & 5d $\to$ 4d &  8{,}301 \\
        2 $\to$ 3 & 10{,}237 & 4d $\to$ 3d &  2{,}246 \\
        3 $\to$ 4 & 16{,}493 & 3d $\to$ 2d & 11{,}451 \\
        4 $\to$ 5 & 12{,}217 & 2d $\to$ 1d &  6{,}815 \\
        5 $\to$ 6 &  8{,}420 & 1d $\to$ 0 &  4{,}907 \\
        \midrule
        \textbf{Total Carga} & \textbf{54{,}051} & \textbf{Total Descarga} & \textbf{40{,}234} \\
        \bottomrule
    \end{tabular}
\end{table}

Restando el área de descarga a la energía absorbida en la carga:

\begin{equation*}
    A \approx 54{,}051 - 40{,}234 = \qty{13{,}817}{kJ/m^3}
\end{equation*}

\begin{equation}
    \boxed{A = \qty{13{,}817}{kJ/m^3} \approx \qty{13817}{J/m^3}}
    \label{eq:resultado_area}
\end{equation}

Este valor representa la energía mecánica disipada internamente por
unidad de volumen de la liga, evidenciando el efecto termodinámico 
irreversible producto de la fricción de sus cadenas poliméricas.